%! AP Statistics — Unit 2, Lesson 2.2 — Student Packet Cover Sheet
\documentclass[10pt]{article}
\usepackage{apstats-article}
\usepackage{apstats-boxes}

%=============================================================================
\begin{document}

% ── Full-bleed navy banner ────────────────────────────────────────────────────
\begin{tikzpicture}[remember picture, overlay]
  \fill[navy] (current page.north west)
    rectangle ([yshift=-0.9in]current page.north east);
\end{tikzpicture}
\vspace{-0.8in}

\begin{center}
  {\color{white}\textbf{\LARGE AP Statistics}}\\[4pt]
  {\color{skymid}\large\textbf{Unit 2: Exploring Two-Variable Data}}\\[6pt]
  {\color{white}\textbf{\large Lesson 2.2 \quad Representing Two Categorical Variables}}
\end{center}

\vspace{0.22in}

\namedateperiod

\vspace{0.18in}

% ── LEARNING TARGETS ─────────────────────────────────────────────────────────
\begin{learningtargetbox}
\small
\begin{enumerate}[leftmargin=1.5em, itemsep=5pt]
  \item \ldots{} construct a two-way (contingency) table and identify joint and marginal
        frequencies.
  \item \ldots{} compute joint and conditional relative frequencies from a two-way table
        and explain the difference between them.
  \item \ldots{} construct and interpret a segmented bar chart using conditional relative
        frequencies to determine whether two categorical variables are associated.
\end{enumerate}
\end{learningtargetbox}

\vspace{0.14in}

% ── TABLE OF CONTENTS ────────────────────────────────────────────────────────
\begin{tocbox}
\small
\renewcommand{\arraystretch}{1.5}
\begin{tabularx}{\linewidth}{@{} c l X r @{}}
\rowcolor{sky}
\textbf{\#} & \textbf{Component} & \textbf{Description} & \textbf{Score} \\
\midrule
1 & Warm-Up       & Spiral review --- relative frequency tables      & \blank{1.2cm} \\
2 & Guided Notes  & Two-way tables, joint/conditional frequencies, segmented bars & \blank{1.2cm} \\
3 & Group Activity & Build and interpret two-way tables by tier        & \blank{1.2cm} \\
4 & Exit Ticket   & Compute conditional frequencies; describe association & \blank{1.2cm} \\
5 & Homework      & Real-data two-way tables + segmented bar charts    & \blank{1.2cm} \\
\midrule
  & \multicolumn{2}{r}{\textbf{Total}} & \blank{1.2cm} \\
\end{tabularx}
\end{tocbox}

\vspace{0.14in}

% ── KEEP IN MIND ─────────────────────────────────────────────────────────────
\begin{remindbox}
\small
Two categorical variables call for a two-way table. Here are the key ideas to keep in mind.

\vspace{6pt}
\begin{tabularx}{\linewidth}{@{} >{\centering\arraybackslash\columncolor{goldbg}}p{0.06\linewidth}
                                    X @{}}
\textbf{\large!} &
\textbf{Joint vs.\ conditional.}\enspace
A joint relative frequency divides a cell by the \emph{grand} total.
A conditional relative frequency divides by a \emph{row or column} total.
To detect association, use conditional frequencies. \\[6pt]
\textbf{\large!} &
\textbf{Condition on the explanatory variable.}\enspace
Identify which variable is explanatory, then divide by its row (or column) total.
Comparing those conditional distributions reveals association. \\[6pt]
\textbf{\large!} &
\textbf{Segmented bars add to 100\%.}\enspace
Each bar represents one group; segments are proportional to conditional relative
frequencies. If the bars look the same, there is no association. \\[6pt]
\textbf{\large!} &
\textbf{Association $\neq$ causation.}\enspace
Even a clear association in a two-way table does not mean one variable causes the
other --- remember Lesson 2.1. \\
\end{tabularx}
\end{remindbox}

\end{document}
